% Drop-in for the delta Cep part of the SPICE paper (mnras_spice_draft_v2.tex), 2026-11-02.
% Replaces \subsubsection{Joint photometric and spectroscopic fit} and its first three paragraphs
% (data/fit set-up, kernel surrogate, results incl. the fitted p-factor), and the three figures
% fig:cepheid_lightcurve, fig:cepheid_spectra, fig:cepheid_teff. The synthetic p-factor paragraphs
% ("The p-factor --- the proportionality constant ...", "The resulting comparison ...") and
% tab:pfactor_line / fig:radvel_ccf are UNCHANGED and stay where they are.
% Label sec:cepheid_fit is kept (referenced from Sec. pulsating_star and Appendix F); new label
% tab:dcep_models.
%
% Also replace in \subsection{Pulsating star} (paragraph starting "To demonstrate"):
%   - "[Fe/H] = +0.06, with [alpha/Fe] and CNO/r/s abundances set to solar"
%       -> "[Fe/H] = +0.08, [C/Fe] = -0.13, [N/Fe] = +0.35 and [O/Fe] = +0.01 \citep{luck2018}"
%   - the sentence "in the joint fit of Section~\ref{sec:cepheid_fit} it is instead held at 2 km/s"
%       -> "and this curve is used in Section~\ref{sec:cepheid_fit}"
%   - the sentences from "The emulator varies smoothly with $\xi$ ..." to "... so $\xi$ is pinned to a node
%     where atmosphere and synthesis agree." -> the paragraph marked XI-REPLACEMENT at the end of this file
%     (the training grid now interpolates the MARCS atmospheres at fixed xi = 2 km/s, so there is no step).
%   - "with $v_{\rm rot} = 12$\,km/s" -> "without rotation \citep{luck2018}"
%   - "as the starting point of the fit in Section~\ref{sec:cepheid_fit}" -> delete
%
% Numbers: Gadi /scratch/y89/mj8805/litrepro/pure/pure_summary.txt (jobs dcep-pure-A/B: dcep_pure_{trahin,luck}_pure_p3,
% dcep_pure_trahin_cont_p3) and cepheid/dcep_vt_level.py. Bib: dcep_literature_section.bib (trahin2021, gordon2023,
% esa1997, anderson2015).
% Figures: python cepheid/dcep_literature_figures.py --trahin out/dcep_pure_trahin_pure_p3.pkl -> paper_plots/cepheid/
%   dcep_photometry.pdf, dcep_spectra.pdf, dcep_teff_vmacro.pdf (copy to figures/).

\subsubsection{The literature model against photometry and spectra}
\label{sec:cepheid_fit}

To test whether one \spice\ model reproduces the photometric and spectroscopic variability of
$\delta$~Cephei at once, we evaluate it at published parameters and compare it with both data sets; nothing
is fitted. The photometry is the Tycho $V_T$ light curve \citep[62 epochs;][]{esa1997}, the only passband of
the \spips\ compilation whose response lies entirely inside the emulator's 3780--6910\,\AA\ range ($V$,
$B$, $H_p$ and $B_T$ extend beyond it). The spectra are the 106 HARPS-N \citep{harpsn} spectra of
\citet{nardetto2017harps}, taken in 2015 and retrieved from the public TNG archive, continuum-normalised
with \texttt{SuppNet} \citep{suppnet} and restricted to twelve 40\,\AA\ windows between 4480 and 6740\,\AA.
(The archive also holds 46 spectra from 2019, whose velocities sit 3.3\,km\,s$^{-1}$ away from the
published orbit described below; we leave them out.)

All inputs are taken from the literature. The radius curve is the \spips\ solution \citep{spips}, scaled to
the distance of \citet{trahin2021} (279.5\,pc, mean radius $43.8\,R_\odot$) at the interferometric angular
diameter, and imposed as the radial $l=0$ displacement. Scaling the radius curve with the distance scales its
pulsation amplitude too, so the model's projection factor is $1.2712 \times 279.5/274.0 = 1.297$, close to
the $p = 1.317$ that \citet{trahin2021} derive at this distance. The mesh is evaluated at the exact phase of
every exposure, so each epoch carries its own velocity field and projected area. For
$T_\mathrm{eff}(\varphi)$ we use two published curves: the \spips\ curve of \citet{trahin2021} (mean 5908\,K)
and a fourth-order Fourier series through the 19 spectroscopic temperatures of \citet{luck2018}. The surface
gravity, microturbulence (2.5--4.0\,km\,s$^{-1}$) and Gaussian line broadening follow the phase-resolved
values of \citet{luck2018}. His $\log g$ comes from the Fe\,\textsc{i}/Fe\,\textsc{ii} ionisation balance
together with his temperatures and microturbulence (mean 1.97, against 1.88 for $5.26\,M_\odot$ at the mean
radius), and his broadening is the total Gaussian width needed to match his own spectra; we read it as a
$1/e$ half-width and remove the 5\,km\,s$^{-1}$ FWHM of his $R = 60\,000$ spectrograph in quadrature (FWHM
9.0--22.7\,km\,s$^{-1}$ over the cycle), since the HARPS-N profile is applied separately. The abundances
also follow \citet{luck2018}, $[\mathrm{Fe/H}] = +0.08$, $[\mathrm{C/Fe}] = -0.13$, $[\mathrm{N/Fe}] = +0.35$
and $[\mathrm{O/Fe}] = +0.01$, and, as in his analysis, there is no rotation. The reddening is
$E(B-V) = 0.092$ \citep{trahin2021} with the $R_V = 3.1$ extinction curve of \citet{gordon2023}. The
systemic velocity is that of the orbit of $\delta$~Cep about its companion \citep{anderson2015} at the 2015
epochs, $-17.11$\,km\,s$^{-1}$. Nothing is adjusted to the data: there is no photometric zero point (the
model magnitudes follow from the radius, temperature, distance and reddening alone), no velocity offset and
no correction of the spectroscopic continuum. Every spectrum and magnitude is the exact per-element
synthesis of Section~\ref{sec:integrating_flux_from_mesh}, with no velocity kernels and no line-correction
template.

Table~\ref{tab:dcep_models} summarises the comparison. With the \citet{trahin2021} temperature curve the
predicted $V_T$ magnitudes are 25\,mmag fainter than observed on average (error-weighted), about the
catalogue error of a single point (25\,mmag) and well within the uncertainty of the reddening alone
($\pm0.021$ in $E(B-V)$ is $\pm70$\,mmag in $V_T$); about that mean the light curve is reproduced to
$\chi^2/N = 1.2$ (Fig.~\ref{fig:cepheid_lightcurve}), and $\chi^2/N = 2.2$ including the offset. The offset
is not uniform over the cycle: around maximum light the residuals scatter about zero, whereas near minimum
light ($\varphi \approx 0.55$--0.75) the model is 50--100\,mmag too faint, which at fixed radius corresponds
to temperatures 50--100\,K too low on that part of the curve. The spectra are reproduced to
$\chi^2/N = 7.1$ against the normalised data as they are, and to 6.4 if a linear continuum correction is
allowed per window and epoch (Fig.~\ref{fig:cepheid_spectra}). The \citet{luck2018} temperatures fit the
spectra equally well (7.0) but the light curve less well ($\chi^2/N = 2.6$, and 2.2 about its mean),
because a Fourier series through 19 points rises too early on the ascending branch
(Fig.~\ref{fig:cepheid_teff}); both curves agree with the independent line-depth-ratio temperatures of
\citet{andrievsky2005} to 24\,K rms, at mean offsets of $+3$ and $+8$\,K. The spectral mismatch is not
uniform in wavelength: $\chi^2/N$ is 8.2--12.1 in the five windows below 5050\,\AA\ and 1.6--2.6 in the four
above 5400\,\AA, for either temperature curve, so it reflects the line data and the emulator in the crowded
blue rather than the pulsation model. As for TZ~Fornacis (Section~\ref{sec:binary_system}), lowering
$[\mathrm{Fe/H}]$ by 0.2\,dex, the offset of our grid's metallicity scale on benchmark stars, lowers the
spectral $\chi^2/N$ by 0.6 at a small cost in the light curve. The line positions test the pulsation
velocity: the cross-correlation velocities of the Mg\,b window follow the model's with a slope of 1.07,
i.e. the observed line shifts call for a velocity field 7 per cent larger than that of $p = 1.297$
($p \approx 1.39$), a projection-factor discrepancy of the size found between line lists and velocity
estimators \citep{nardetto2017harps}. The radius curve, the temperatures, the line broadening and the
binary orbit of four independent analyses thus combine, through \spice, into a model that reproduces the
light curve and the line profiles over the whole cycle without any quantity adjusted to these data.

\begin{table}
    \centering
    \caption{$\delta$~Cephei at literature parameters, with nothing fitted, against the Tycho $V_T$ light curve
    (62 points, catalogue errors; no zero point) and 106 HARPS-N spectra in twelve windows. $\Delta V_T$ is
    the mean observed minus model magnitude; the spectral $\chi^2/N$ in parentheses allows a linear continuum
    correction per window and epoch.}
    \begin{tabular}{lccc}
        \hline
        $T_\mathrm{eff}(\varphi)$ & $\chi^2_{V_T}/N$ & $\Delta V_T$ [mmag] & $\chi^2_{\rm spec}/N$ \\
        \hline
        \citet{trahin2021} & 2.17 & $-25$ & 7.09 (6.38) \\
        \citet{luck2018}   & 2.63 & $-16$ & 6.96 \\
        \hline
    \end{tabular}
    \label{tab:dcep_models}
\end{table}

\begin{figure}
    \centering
    \includegraphics[width=0.9\columnwidth]{figures/dcep_photometry.pdf}
    \caption{Tycho $V_T$ light curve of $\delta$~Cephei against the \spice\ model at literature parameters
    (Section~\ref{sec:cepheid_fit}; \citealt{trahin2021} temperature curve), with the residuals beneath; the
    dashed line is their weighted mean. The model is the exact per-element synthesis; no zero point is
    adjusted.}
    \label{fig:cepheid_lightcurve}
\end{figure}

\begin{figure*}
    \centering
    \includegraphics[width=\textwidth]{figures/dcep_spectra.pdf}
    \caption{HARPS-N spectra of $\delta$~Cephei \citep{nardetto2017harps} at three pulsation phases (marked
    on the $V_T$ light curve, left) against the per-element synthesis at literature parameters, in the
    4480--4520\,\AA\ window (the bluest and least well reproduced of the twelve), with residuals beneath. Each epoch's velocity field follows
    from the mesh at its exact phase.}
    \label{fig:cepheid_spectra}
\end{figure*}

\begin{figure}
    \centering
    \includegraphics[width=\columnwidth]{figures/dcep_teff_vmacro.pdf}
    \caption{Top: the two temperature curves of the literature models, \citet{trahin2021} (solid) and the
    Fourier series through the spectroscopic temperatures of \citet{luck2018} (dashed, points), with the
    line-depth-ratio temperatures of \citet{andrievsky2005}; ticks mark the phases of the 106 HARPS-N epochs.
    Bottom: the line broadening (FWHM) of \citet{luck2018} used in both models, with his instrumental profile
    removed; dots mark the HARPS-N epochs.}
    \label{fig:cepheid_teff}
\end{figure}

% ---------------------------------------------------------------------------
% XI-REPLACEMENT (for \subsection{Pulsating star}):
%
% The training spectra of the emulator are computed on model atmospheres interpolated in
% $T_\mathrm{eff}$, $\log g$ and $[\mathrm{Fe/H}]$ at a fixed atmospheric microturbulence of
% 2\,km\,s$^{-1}$ (Appendix~\ref{sec:appendix_spectra_emulators}), with $\xi$ entering only the line
% formation, so the emulated flux varies smoothly with $\xi$; the phase-dependent $\xi$ of
% \citet{luck2018} can therefore be used directly.
% ---------------------------------------------------------------------------
